

\begin{frame}{Example - cable robot, 1}
	% \framesubtitle{General form}
	\begin{flushleft}
		
		Consider a cable-driven parallel robot.
		
		% TODO: \usepackage{graphicx} required
		\begin{figure}
			\centering
			\includegraphics[width=0.5\linewidth]{"cable robot"}
		\end{figure}
		
		
		The points where the cables are attached are $O_1$, ..., $O_n$ with coordinates $\bo{r}_1$, ..., $\bo{r}_n$. The center of mass of the robot has coordinates $\bo{r}_0$. The directions of the cables are given as $\bo{p}_1$, ..., $\bo{p}_n$ and the gravity force is $\bo{f}_0$. The mass of the robot is $m$, inertia - $\bo{I}$.
		
		\bigskip
		
		Find what forces the cables need to apply to achive acceleration as close to $\bo{a}^*$ as possible, without angular acceleration.
		
	\end{flushleft}
\end{frame}



\begin{frame}{Example - cable robot, 2}
	% \framesubtitle{General form}
	\begin{flushleft}
		
		We define the force acting from the $i$-th cable as $ \bo{f}_i  = \bo{p}_i \alpha_i$, where $\alpha_i$ is the magnitude of the force (assuming $||\bo{p}_i|| = 1$). We can write newton equation:
		
		\begin{align}
			m \bo{a} = \bo{f}_0 + \sum_{i = 1}^{n} \bo{p}_i \alpha_i
		\end{align}
		
		Similar, we write Euler equation:
		
		\begin{align}
			\bo{I}\varepsilon = [\bo{r}_0]_\times \bo{f}_0 + \sum_{i = 1}^{n} [\bo{r}_i]_\times \bo{p}_i \alpha_i
		\end{align}
		
		where $[\bo{r}]_\times$ is a skew-symmetric representation of a vector; $[\bo{r}]_\times \bo{x} = \bo{r} \times \bo{x}$.
		
	\end{flushleft}
\end{frame}


\begin{frame}{Example - cable robot, 3}
	% \framesubtitle{General form}
	\begin{flushleft}
		
		Thus, we the problem takes form:
		
		%
		\begin{equation}
			\begin{aligned}
				& \underset{\alpha_i, \bo{a}, \varepsilon}{\text{minimize}}
				& & (\bo{a}^* - \bo{a})\T (\bo{a}^* - \bo{a}) + \varepsilon\T \varepsilon, 
				\\
				& \text{subject to}
				& & \begin{cases}
					m \bo{a} = \bo{f}_0 + \sum_{i = 1}^{n} \bo{p}_i \alpha_i,
					\\
					\bo{I} \varepsilon = [\bo{r}_0]_\times \bo{f}_0 + \sum_{i = 1}^{n} [\bo{r}_i]_\times \bo{p}_i \alpha_i.
				\end{cases}
			\end{aligned}
		\end{equation}
		
		\bigskip
		
		After solving the problem, we find the forces:
		
		\begin{align}
			\bo{f}_i  = \bo{p}_i \alpha_i
		\end{align}
		
	\end{flushleft}
\end{frame}